\documentclass[a4paper]{article}
\usepackage{luatexja}
\usepackage{amsmath}
\usepackage{amssymb}
\usepackage{geometry}
\usepackage{hyperref}

\geometry{margin=25mm}

\newcommand{\dstrec}{\textbf{記録}}
\newcommand{\dstwork}{\textbf{作業模型}}
\newcommand{\dstopen}{\textbf{未定義}}

\title{二重時空の多体計算\\プログラマのための実装仕様}
\author{https://github.com/hypernumbernet}
\date{\today}

\begin{document}

\maketitle

\begin{abstract}
本書は、二重時空理論（DST）の機械検査済み定義から、多体計算プログラムの仕様を切り出す。正本は Lean~4 の開発である。論文草稿と Lean が食い違う箇所では Lean を採る。各アルゴリズムには地位を付ける。\dstrec{}は Lean が定義または証明している式である。\dstwork{}は書かれた作用の近傍に置かれた別ラグランジアンである。\dstopen{}は多体の時間結合のように、呼び出し側が宣言しない限りプログラムに存在しない界面である。内在角の時間発展は閉形式で返す。層の力は有限集合上の静的な和である。チャート上のシュヴァルツシルト比較と、粒子内在角の積分は別トラックである。
\end{abstract}

\tableofcontents

\section{読み方}
\label{sec:reading}

本書の読者は、DST の多体計算を書くプログラマである。物理の総説をここで再構成しない。実装が合わせる相手は、リポジトリの Lean モジュールである。モジュール文書と定理が仕様であり、論文の段落は索引である。

\subsection{三つの地位}

関数を置く前に、その関数の地位をファイル先頭のコメントへ書く。

\begin{description}
\item[\dstrec] Lean にある定義または定理と、引数の意味まで一致させる。名前は Lean の識別子をそのまま使う。
\item[\dstwork] 動態の論文が「書かれた作用ではない」と分けている近傍系である。識別子に \texttt{working} を付ける。記録用のソルバから呼ばない。
\item[\dstopen] Lean に積分器も接触則も無い界面である。プログラムが値を埋めるなら、呼び出し側が渡した仮定としてログへ残す。論文の運動方程式へ昇格させない。
\end{description}

地位の違う式を一つの更新関数に混ぜると、テストがどの主張を見ているか分からなくなる。モジュール境界は第\ref{sec:modules}章で固定する。

\subsection{二つのトラック}

DST には、少なくとも次の二つの計算がある。

\begin{enumerate}
\item 粒子の内在角 \((\phi_a,\theta_a)\) の時間発展。書かれた作用の閉形式は第\ref{sec:written}章である。
\item 共有テトラッドのチャート上の外部軌道。水星の近日点余剰は、外部の軌道要素を動径回帰へ入れた量である。太陽光線の偏折は、外部のヌル軌道の方位角余剰である。
\end{enumerate}

チャートの数値と、内在角の閉形式の出力は、別のテストで見る。対応は第\ref{sec:checks}章に書く。層の節に出る \(\alpha,\beta\) は干渉因子 \(\gamma_s\) の引数であり、粒子状態のラピディティとは別の読みである。層力は粒子の内在角を読まない。

\subsection{多体という言葉}

Lean が多体について固定しているのは、有限集合上の静的な力和である。時間ステップは取らない。粒子数を変える反応も置かない。時刻を進める一段が必要なら、第\ref{sec:contract}章の仮定をログに残したうえで、作業模型としてだけ実行する。

\section{一粒子の状態}
\label{sec:state}

地位は\dstrec{}である。型は \texttt{Algebra.Operations.TorsionParams} であり、場は次の二つである。

\begin{verbatim}
struct TorsionParams:
    alpha: array[3] of real   # usual rapidity phi
    beta:  array[3] of real   # dual angle theta
\end{verbatim}

添字 \(a\in\{0,1,2\}\) は Lean の \texttt{Fin 3} である。三軸は、書かれた作用の中では別々に積分する。軸を混ぜる項を、この記録のソルバへ足さない。

\subsection{共通ラピディティと不整合}

一軸ごとに
\[
\sigma_a=\alpha_a+\beta_a,\qquad
\delta_a=\alpha_a-\beta_a
\]
と置く。逆は
\[
\alpha_a=\frac{\sigma_a+\delta_a}{2},\qquad
\beta_a=\frac{\sigma_a-\delta_a}{2}
\]
である。\texttt{DualRotorFlow.usualFrom} と \texttt{dualFrom} は、時刻の関数 \(\sigma,\delta\) にこの一次式を点ごとに適用する。プログラムは \(\alpha,\beta\) と \(\sigma,\delta\) を別々の状態として持たない。\(\sigma(t)\) と \(\delta(t)\) を閉形式で計算し、角はこの一次式で作る。

\subsection{ねじれスカラーと符号なし質量}

実装する式は \texttt{Invariant.J} と \texttt{Invariant.mass} である。
\begin{align}
J(p)&=\frac12\sum_{a=0}^{2}\bigl(\alpha_a^2-\beta_a^2\bigr)
=\frac12\sum_{a=0}^{2}\delta_a\sigma_a,\\
M(p)&=\frac12\sum_{a=0}^{2}\bigl(\alpha_a^2+\beta_a^2\bigr)
=\frac14\sum_{a=0}^{2}\bigl(\sigma_a^2+\delta_a^2\bigr).
\end{align}
生成子のキリング形式を \(16\) で割った値は \(J\) ではない。恒等式は
\[
J(p)=4\cdot\Bigl(\frac1{16}\,\texttt{omegaTorsionGeneratorKilling}(p)\Bigr)
\]
であり、右辺の括弧は \(\frac18\sum_a(\alpha_a^2-\beta_a^2)\) である。係数 \(8\sum_a(\alpha_a^2-\beta_a^2)\) も \(J\) ではない。

正規化は
\[
J_{\mathrm{norm}}(p)=\frac{8}{3\pi^2}J(p)
\]
であり、識別子は \texttt{JNormalized} である。係数は正なので、\(|J_{\mathrm{norm}}|\le 1\) と \(|J|\le 3\pi^2/8\) はすべての状態で同値である。許容錐の上では、この不等式が成り立つ。定理は \texttt{torsion\_bound\_continuous} と \texttt{torsion\_bound\_raw\_continuous\_pi\_sq} である。錐の上で \(|J_{\mathrm{norm}}|=1\) となるのは、全軸が純双曲か全軸が純楕円かのときである。

\(J(p)=0\) は \(\sum_a\alpha_a^2=\sum_a\beta_a^2\) と同値である。各軸で \(\delta_a=0\) となることは、この等式の特別な場合である。軸をまたいで平方が打ち消し合っても \(J\) は消える。

一軸だけが非零のとき、\(J=\frac12\delta\sigma\)、\(M=\frac14(\sigma^2+\delta^2)\) はその軸への制限である。書かれた作用の剛性ポテンシャルは別の量で、
\[
\texttt{stiffnessPotential}(m,p)=\frac m2\sum_a(\alpha_a-\beta_a)^2
\]
である。これを \(J\) へ代入しない。一軸に限れば、\(m\neq 0\) のとき両者が \(m\) 倍で一致するのは \(\alpha=\beta\) または \(\beta=0\) のときに限る。

\(J=0\) は相対ローターが単位元であることを意味しない。一軸で \(\delta=0\) かつ \(\sigma\neq 0\) なら二つの角は等しく、\(\alpha=\beta=\sigma/2\neq 0\)、\(J=0\)、\(M=\sigma^2/4>0\) である。このとき相対ローターは単位元ではない。証人は \texttt{axisParams}\(1\,1\) であり、定理は \texttt{RelativeRotor.J\_zero\_not\_relativeRotor\_one} である。

\subsection{許容錐}

\texttt{Admissible.IsAdmissibleContinuous} は、各軸で
\[
\alpha_a\ge 0,\qquad \beta_a\ge 0,\qquad \alpha_a+\beta_a\le\frac\pi2
\]
が成り立つことである。\(\sigma,\delta\) へ訳すと、各軸で
\[
\sigma_a\in\Bigl[0,\frac\pi2\Bigr],\qquad |\delta_a|\le\sigma_a
\]
である。主枝の絶対値条件 \(|\alpha_a+\beta_a|\le\pi/2\) だけでは \(J\) は抑えられない。検査関数は連続許容の三つ組を見る。

許容錐を外れた状態の扱いは第\ref{sec:numerics}章で拒否とする。錐の中へ射影する関数は記録に無い。

\subsection{一つの質量にまとめない}

\texttt{RestMass} は、次の三つを一つの数にしない。

\begin{enumerate}
\item 剛性ポテンシャル \texttt{stiffnessPotential}。係数 \(m\) は入力である。式は前節のとおりである。
\item 符号なし質量 \(M\)。ラピディティのユークリッド二次形式である。
\item 自由落下。\(F=ma\) という識別は主張されていない。加速度の上限も主張されていない。
\end{enumerate}

段の更新に慣性パラメータが要るなら、呼び出し側が別に渡す。剛性 \(m\) でも \(M\) でもない。詳細は第\ref{sec:contract}章である。

\section{位置は二種類ある}
\label{sec:position}

地位は\dstrec{}である。多体の層力と、関係としての位置は、別の型に入れる。

\subsection{チャート上の点}

\texttt{LayerForce} の位置は、実内積空間 \(E\) の点である。三次元の計算では \(E=\mathbb{R}^3\) で足りる。粒子の識別子を \(\iota\) とすると、配置は
\begin{verbatim}
pos:  map iota -> Vec3
\end{verbatim}
である。分離は \(\lVert\mathrm{pos}(i)-\mathrm{pos}(j)\rVert\) であり、このノルムが層の引数 \(r\) になる。この \(E\) は粒子の内在角が住む代数ではない。

\subsection{三球上の関係}

\texttt{RelationalPosition.S3} は単位四元数である。
\begin{verbatim}
struct S3:
    val: quaternion
    # normSq(val) = 1
\end{verbatim}
\(b\) の \(a\) に対する位置は相対ローター
\[
\rho(a,b)=a^{-1}b
\]
であり、識別子は \texttt{relative} である。\(\rho(a,b)=1\) は \(a=b\) と同値である。左からの同時移動は \(\rho\) を変えない。

この \(\rho\) は、粒子内部の不整合ローター \(\Omega\) ではない。チャート座標の位置演算子でもない。\(\rho\) を \(\mathbb{R}^3\) の差と見なして、層力の \(r\) に渡さない。\(S^3\) 上の測地線を、ニュートン軌道の代わりに積分しない。

\subsection{状態を分ける}

最小の実行時状態は次の三層に割る。

\begin{verbatim}
struct ChartBody:
    id: iota
    pos: Vec3
    mass_input: real    # factor in G * m_i * m_j
    charge_input: real  # factor in k * q_i * q_j

struct WrittenSeed:
    axis: int           # 0, 1, or 2
    sigma0: real
    sigmadot0: real
    delta0: real
    nu: real            # constant delta-dot
    stiffness: real     # input m

struct RelationalSample:
    element: S3         # not read by layer_force
\end{verbatim}

\texttt{ChartBody} は第\ref{sec:force}章だけが読む。\texttt{WrittenSeed} は第\ref{sec:written}章だけが読む。\texttt{RelationalSample} は関係の検査に残し、時間ループの位置更新へ接続しない。一つの構造体に三層の場を並べることはできる。そのとき更新関数の引数リストで、読む場を制限する。

\section{入力であって導出でない量}
\label{sec:inputs}

地位は\dstopen{}を、入力スロットとして固定する話である。次の量を幾何から計算する関数は置かない。

\begin{itemize}
\item 重力定数 \(G\)。
\item 層の長さ \(\ell\)。チャネルが二つあるときは \(\ell_G\) と \(\ell_E\) を別引数にする。
\item 書かれた作用の剛性 \(m\)。
\item 電荷、およびクーロン前置 \(k\)。
\item 微細構造定数、ボーア半径、核の長さ \(r_1\)。
\item 銀河スケールの \(a_0\)。
\end{itemize}

\texttt{Gravity.SI} の有理数は、CODATA 型の外部入力を回帰で使うためのスタンドインである。代数から導出された値ではない。本番のソルバがこの有理数を既定値として焼き込まない。テストだけが、回帰検査の数値範囲と比べるために明示して渡す。

\subsection{受け取り方}

定数は有理数で保持できる。浮動小数へ落とすなら、落とした桁と、比較に使う区間を同じレコードに書く。
\begin{verbatim}
struct Inputs:
    G: rational
    ell_G: rational
    ell_E: rational
    k_coulomb: rational
    stiffness: rational
    note: string   # which external table, which digits
\end{verbatim}
ゼロ除算と符号の検査は、この値を力の式へ入れる前に行う。\(G\) の符号、\(\ell\) の符号、剛性の符号は、呼び出し側の責任としてログへ残す。Lean の層力定理は、正の \(\kappa=G m_i m_j\) と正の \(r\) を仮定して符号を述べている。

\subsection{単位}

内在角 \(\alpha,\beta,\sigma,\delta\) は無次元である。\(J\) と \(M\) も無次元である。チャート上の \(r\) は \(\ell\) と同じ単位で測る。\(\ell\) を fm で渡した力と、\(\ell\) をメートルで渡した力を、一つの和に入れない。SI へ変換する関数が要るなら、変換係数も \texttt{Inputs} のフィールドにし、幾何関数の中へ書かない。

\section{書かれた一粒子作用}
\label{sec:written}

地位は\dstrec{}である。ソースは \texttt{DualRotorDynamics} と \texttt{DualRotorFlow} である。実装は数値 ODE を先に書かない。この章の関数は閉形式を返す。

\subsection{密度とオイラー--ラグランジュ方程式}

Lean が定義している密度は一軸のものである。
\[
L=\frac12\dot\phi^2-\frac12\dot\theta^2-\frac m2(\phi-\theta)^2
\]
識別子は \texttt{paperLagrangian} である。この密度に軸をまたぐ項は無い。三軸を同時に進めるときは、同じ密度を軸ごとの初期値で別々に評価する。

正しいオイラー--ラグランジュ方程式 \texttt{PaperActualEL} は
\[
\ddot\phi=-m(\phi-\theta),\qquad
\ddot\theta=-m(\phi-\theta)
\]
である。したがって
\[
\ddot\delta=0,\qquad
\ddot\sigma=-2m\delta.
\]
書かれた方程式は \(\ddot\delta+m\delta=0\) ではない。定理 \texttt{paperActualEL\_not\_oscillator} の証人は
\[
(m,\phi,\theta,\ddot\phi,\ddot\theta)=(1,1,0,-1,-1)
\]
であり、この点では \(\ddot\delta=0\) かつ \(-m\delta=-1\) である。第\ref{sec:working}章の同符号模型はさらに別の方程式 \(\ddot\delta+2m\delta=0\) であり、ここへ混ぜない。

\subsection{閉形式}

初期値を \(\delta(0)=\delta_0\)、\(\dot\delta=\nu\)（定数）、\(\sigma(0)=\sigma_0\)、\(\dot\sigma(0)=\dot\sigma_0\) とする。解は
\begin{align}
\delta(t)&=\delta_0+\nu t,\\
\dot\sigma(t)&=\dot\sigma_0-2m\delta_0 t-m\nu t^2,\\
\sigma(t)&=\sigma_0+\dot\sigma_0 t-m\delta_0 t^2-\frac{m\nu}{3}t^3
\end{align}
である。識別子は \texttt{writtenDelta}、\texttt{cubicRapiditySpeed}、\texttt{writtenSigma} である。角へ戻すときは
\begin{verbatim}
phi(t) = (sigma(t) + delta(t)) / 2
theta(t) = (sigma(t) - delta(t)) / 2
\end{verbatim}
とする。\(\ddot\phi=\ddot\theta=-m\delta(t)\) が各時刻で成り立つことを、実装後の検査に入れる。

\subsection{運動量とヤコビ積分}

共役運動量に剛性 \(m\) は掛からない。
\begin{align}
p_\sigma&=\frac12\dot\delta=\frac12(\dot\phi-\dot\theta),\\
p_\delta&=\frac12\dot\sigma=\frac12(\dot\phi+\dot\theta).
\end{align}
識別子は \texttt{paperMomentumSigma} と \texttt{paperMomentumDelta} である。書かれた流れの上では \(p_\sigma\) は保存し、\(\dot p_\delta=-m\delta\) である。

ヤコビ積分は
\[
E=\frac12\dot\phi^2-\frac12\dot\theta^2+\frac m2\delta^2
=\frac12\dot\sigma\dot\delta+\frac m2\delta^2
\]
であり、識別子は \texttt{paperEnergy} である。\texttt{PaperActualEL} の上で \(\dot E=0\) である。閉形式の上では、すべての時刻で
\[
E(t)=\frac12\nu\dot\sigma_0+\frac m2\delta_0^2
\]
である。定理は \texttt{written\_flow\_energy} である。\(E\) は不定であり、正にも負にもなりうる。この保存は、振動数の存在を意味しない。

\subsection{呼び出し境界}

\begin{verbatim}
function written_at(seed, t) -> (phi, theta, phidot, thetadot):
    delta = seed.delta0 + seed.nu * t
    sigmadot = seed.sigmadot0
             - 2 * seed.stiffness * seed.delta0 * t
             - seed.stiffness * seed.nu * t * t
    sigma = seed.sigma0
          + seed.sigmadot0 * t
          - seed.stiffness * seed.delta0 * t * t
          - (seed.stiffness * seed.nu / 3) * t * t * t
    phi = (sigma + delta) / 2
    theta = (sigma - delta) / 2
    phidot = (sigmadot + seed.nu) / 2
    thetadot = (sigmadot - seed.nu) / 2
    return (phi, theta, phidot, thetadot)
\end{verbatim}

この関数は力も位置も受け取らない。駆動 \(u\)、振動子、暴走、任意ポテンシャルは第\ref{sec:working}章の別モジュールである。

\section{近傍模型を隔離する}
\label{sec:working}

地位は\dstwork{}である。ソースは \texttt{DualRotorDynamics} の近傍系、\texttt{DualRotorFlow} の双曲線解と調和解、\texttt{DualRotorRigidity} の傾きジェットである。ファイル名と関数名に \texttt{working} を付ける。第\ref{sec:written}章のモジュールを import して式を再利用してよい。逆方向、つまり作業模型を記録ソルバの既定経路にする import は禁止する。

三つの二次系は隣り合っている。運動項の符号が、\(\delta\) を力学座標にするかどうかを決める。書かれた作用は反対符号であり、振動しない。振動子は同符号であり、書かれた作用ではない。残る系は暴走である。

\subsection{暴走}

\texttt{PaperClaimedEL} は
\[
\ddot\phi=m\delta,\qquad \ddot\theta=-m\delta
\]
である。チャネルでは
\[
\ddot\delta=2m\delta,\qquad \ddot\sigma=0.
\]
\(\kappa=\sqrt{2m}\)（\(m\ge 0\)）として、閉形式は
\begin{align}
\delta(t)&=A\cosh(\kappa t)+B\sinh(\kappa t),\\
\sigma(t)&=\sigma_0+\dot\sigma_0 t
\end{align}
である。識別子は \texttt{runawayDelta} と \texttt{runawaySigma} である。この解を \texttt{written\_at} の戻り値にしない。

\subsection{振動子}

同符号の密度は
\[
L_{\mathrm{osc}}=\frac12\dot\phi^2+\frac12\dot\theta^2-\frac m2\delta^2
\]
であり、識別子は \texttt{oscillatorLagrangian} である。方程式 \texttt{OscillatorEL} は
\[
\ddot\delta+2m\delta=0,\qquad \ddot\sigma=0
\]
である。角振動数は \(\sqrt{2m}\) である。単位 \(\hbar=c=1\) で、この値がコンプトン振動数 \(m\) と一致するのは \(m=0\) と \(m=2\) に限る。定理は \texttt{sqrt\_two\_mul\_eq\_compton\_iff} である。

閉形式は
\[
\delta(t)=A\cos\bigl(\sqrt{2m}\,t\bigr)+B\sin\bigl(\sqrt{2m}\,t\bigr),
\quad
\sigma(t)=\sigma_0+\dot\sigma_0 t
\]
である。識別子は \texttt{harmonicMismatch} と \texttt{inertialRapidity} である。ヤコビ積分
\[
E_{\mathrm{osc}}=\frac12\dot\phi^2+\frac12\dot\theta^2+\frac m2\delta^2
\]
は \(m\ge 0\) のとき非負であり、\(\delta\) の運動を束縛する。書かれた作用の \(E\) とは符号も意味も違う。

\subsection{双対チャネルの加速度}

書かれた密度に、双対側の加速度 \(u\) を足したジェットが \texttt{PaperSourcedEL} である。
\[
\ddot\phi=-m\delta,\qquad
\ddot\theta=-m\delta+u,
\]
したがって
\[
\ddot\delta=-u,\qquad
\ddot\sigma=-2m\delta+u.
\]
\(u=0\) のときだけ \texttt{PaperActualEL} に戻る。\(u\) が定数でも、不整合は復元振動数を持たない。二重積分である。

通常側の力 \(v\) も認める \texttt{PaperBothSourcedEL} では
\[
\ddot\phi=-m\delta+v,\qquad
\ddot\theta=-m\delta+u
\]
であり、ヤコビ積分の時間微分は \(v\dot\phi-u\dot\theta\) である。\(\ddot\phi=0\) を課すと \(v=m\delta\) が一意に決まる。この \(v\) は自由解ではない。通常角が瞬間的に静止していれば、拘束は仕事をしない。

同符号模型に \(u\) を足すと
\[
\ddot\delta+2m\delta=-u,\qquad \ddot\sigma=u
\]
である。\(m\neq 0\) の定数駆動は、平衡不整合を \(-u/(2m)\) だけずらす。\(u\) を実験室のファラデー双一次量と同定するのは模型の選択であり、導出ではない。

\subsection{任意の傾き}

\texttt{OppositePotentialEL} は、傾き \texttt{slope} \(=V'(\delta)\) に対して
\[
\ddot\phi=\ddot\theta=-\mathrm{slope}
\]
と置く。したがって \(\ddot\delta=0\) かつ \(\ddot\sigma=-2V'(\delta)\) である。調和剛性では \(\mathrm{slope}=m\delta\) であり、書かれた作用の力と一致する。反対符号のままでは、\(V\) の形を変えても \(\delta\) は自由なままである。零でない遅れは、傾きが消える点に留まらない。

\texttt{SameSignPotentialEL} は
\[
\ddot\phi=-\mathrm{slope},\qquad \ddot\theta=\mathrm{slope}
\]
である。したがって \(\ddot\delta=-2V'(\delta)\) かつ \(\ddot\sigma=0\) である。四次井戸
\[
V(\delta)=\frac\lambda4(\delta^2-v^2)^2,\qquad
V'(\delta)=\lambda\delta(\delta^2-v^2)
\]
の識別子は \texttt{quarticPotential} と \texttt{quarticSlope} である。\(\lambda\neq 0\) のとき傾きが消えるのは \(\delta=0,\pm v\) である。非零の井戸での曲率は \(4\lambda v^2\) であり、識別子は \texttt{radialFreqSq} である。この曲率をコンプトン値へ固定する関数は置かない。ワインバーグ角も \(W,Z\) の質量も、このモジュールから返さない。

\subsection{API}

\begin{verbatim}
module working_models:
    function runaway_at(...)
    function oscillator_at(...)
    function sourced_accelerations(m, delta, u, v)
    function opposite_slope_el(slope)
    function samesign_slope_el(slope)
    function quartic_slope(lam, v, delta)
\end{verbatim}

\texttt{written\_at} は上のどれも呼ばない。テストは、同一の \((m,\phi,\theta)\) に対して \texttt{PaperActualEL} の加速度と \texttt{PaperClaimedEL} の加速度が一致しないサンプルを一つ持つ。

\section{層因子}
\label{sec:layer}

地位は\dstrec{}である。ソースは \texttt{ElectronShell} と \texttt{TorsionalLayer} である。

\subsection{二つの干渉因子}

一般の対は
\begin{align}
\gamma_s(\alpha,\beta)
&=\cosh\frac\alpha2\cos\frac\beta2-\sinh\frac\alpha2\sin\frac\beta2,\\
\gamma_B(\alpha,\beta)
&=-\cosh\frac\alpha2\sin\frac\beta2+\sinh\frac\alpha2\cos\frac\beta2,\\
\gamma_{\mathrm{eff}}(\alpha,\beta)
&=\gamma_s^2+\gamma_B^2
\end{align}
である。\(\gamma_{\mathrm{eff}}\ge 0\) は常に成り立つ。力の符号には使わない。固有時の読みと、層の符号は別の関数である。

等スケールは \(x=\ell/(2r)\) と置き、
\[
\gamma_s^{\mathrm{eq}}(x)=\cosh x\cos x-\sinh x\sin x
\]
とする。識別子は \texttt{gammaSEqual} である。これは \(\gamma_s(2x,2x)\) に等しい。
\[
\gamma_s^{\mathrm{eq}}(0)=1,\qquad
\frac{d}{dx}\gamma_s^{\mathrm{eq}}(x)=-2\cosh x\sin x.
\]
導関数が消えるのは \(\sin x=0\) のときである。

\subsection{第一の正根}

\texttt{resonanceRoot1} は、開区間 \((1/2,1)\) 内で \(\tanh x\tan x=1\) を満たす唯一の点である。その点で \(\gamma_s^{\mathrm{eq}}=0\) である。回帰検査は、より狭い範囲
\[
\frac\pi4<\texttt{resonanceRoot1}<1
\]
を要求する。根の閉形式を定数として埋めない。テストは、二分法または同等の手続きで見つけた根がこの区間に入り、その点で \(\gamma_s^{\mathrm{eq}}\) の絶対値が許容誤差未満であることを見る。

第一斥力殻は、位相が \((\texttt{resonanceRoot1},\pi)\) に入る半径である。そこでは \(\gamma_s^{\mathrm{eq}}<0\) である。層半径は
\[
r=\frac{\ell}{2x}
\]
であり、識別子は \texttt{layerRadius} である。

\subsection{零と近接点}

\(\gamma_s^{\mathrm{eq}}(x)=0\) の半径では、層力の分母が消える。そのペアに対する力の評価は失敗として返し、段を進めない。\(r=0\) も同様に失敗である。消滅する分離は接触則ではない。接触、合体、反射を表す分岐を、この章へ書かない。

切断半径 \(r_{\mathrm{cut}}>0\) を模型が欲しがるなら、\texttt{Inputs} の外から渡す。\(\ell\) から \(r_{\mathrm{cut}}\) を計算しない。\(r<r_{\mathrm{cut}}\) のペアは、力を数値で埋めず、拒否する。

\subsection{実装}

\begin{verbatim}
function gamma_s_equal(x) -> real:
    return cosh(x) * cos(x) - sinh(x) * sin(x)

function phase(ell, r) -> real or reject:
    if r == 0: reject
    return ell / (2 * r)

function layer_gamma(ell, r) -> real or reject:
    g = gamma_s_equal(phase(ell, r))
    if g == 0: reject
    return g
\end{verbatim}

浮動小数では、零の判定を絶対許容 \(\varepsilon_\gamma\) で行う。\(\varepsilon_\gamma\) も入力であり、理論定数ではない。許容の内側で分母を小さい数に置き換えて力を連続化しない。

\section{静的な多体力}
\label{sec:force}

地位は\dstrec{}である。この章の関数は時刻を受け取らない。加速度も返さない。

\subsection{重力だけの動径項}

\(\kappa\) は入力 \(G m_1 m_2\) である。
\[
\texttt{gravityLayerTerm}(\kappa,\ell,r)
=-\frac{\kappa}{\gamma_s^{\mathrm{eq}}(\ell/(2r))\,r^2}.
\]
正の値は分離を増やす。\(\kappa>0\)、\(r>0\)、\(\gamma_s^{\mathrm{eq}}<0\) ならば項は正である。同じ仮定で \(\gamma_s^{\mathrm{eq}}>0\) ならば項は負であり、分離を減らす向きである。

点 \(p\) が点 \(c\) から受けるベクトルは
\[
\texttt{gravityLayerPush}(\kappa,\ell,c,p)
=\frac{\texttt{gravityLayerTerm}(\kappa,\ell,\lVert p-c\rVert)}{\lVert p-c\rVert}\,(p-c)
\]
である。外向き単位分離との内積が、動径項に一致する。反対称性
\[
\texttt{gravityLayerPush}(\kappa,\ell,c,p)
=-\texttt{gravityLayerPush}(\kappa,\ell,p,c)
\]
を、二体テストの等式として使う。

有限集合 \(S\) の粒子 \(i\) に働く力は
\[
\texttt{gravityLayerForce}_i
=\sum_{j\in S\setminus\{i\}}
\texttt{gravityLayerPush}(G m_i m_j,\ell,\mathrm{pos}_j,\mathrm{pos}_i)
\]
である。和は \texttt{erase i} であり、自己項は無い。実装手順は次の順にする。

\begin{enumerate}
\item 二体のスカラー \texttt{gravityLayerTerm} を、引力層と斥力層のサンプルで評価する。
\item 二体のベクトルが反対称であることを見る。
\item 三体で、粒子 \(i\) の和に \(i\) 自身が現れないことを見る。
\item \(N\) 体へ広げる。各項の失敗（\(r=0\) または \(\gamma=0\)）は、和を部分的に返さず、呼び出し全体を拒否する。
\end{enumerate}

\subsection{二チャネル}

識別子 \texttt{pairChannelForce} は分離を増やす向きのスカラーである。
\[
-\frac{\kappa}{\gamma_G r^2}+\frac{\mathrm{coul}}{\gamma_E r^2}.
\]
\(\kappa=G m_1 m_2\)、\(\mathrm{coul}=k q_1 q_2\) である。\(\gamma_G,\gamma_E,r\) が非零のとき、この量が零になることと
\[
\kappa\,\gamma_E=\mathrm{coul}\,\gamma_G
\]
は同値である。二つの干渉因子を固定したままでは、この条件に半径は現れない。正の半径がすべて零か、どれも零でないかである。\(\gamma\) を \(r\) から計算し直した合成関数の零を、平衡半径として探索するループは書かない。

重力項は、電荷項を \(0\) にし \(\gamma_E=1\) とした二チャネル力と一致する。電磁項を \texttt{gravityLayerForce} の和へ混ぜない。二チャネルが要る呼び出しは、別関数 \texttt{pair\_channel\_force} を使う。

\(\kappa>0\) かつ \(\mathrm{coul}<0\) の異種対では、干渉因子が同符号なら分離を増やす力は零にならない。両方の因子が正なら力は負であり、両方とも負なら力は正である。実装はそのサンプルをテストに残す。スペクトルも束縛エネルギーも、この関数の戻り値にしない。

\subsection{クーロン項と点電荷}

スカラー
\[
\texttt{coulombForce}(k,e,\gamma_s,r)=\frac{k e^2}{\gamma_s r^2}
\]
の符号は \(\gamma_s\) の符号である。円軌道の読み \(v^2=ke^2/(m\gamma_s r)\) は、このスカラーと \(r/m\) の積として書かれている。その \(m\) は呼び出し側のパラメータであり、剛性からも \(M\) からも計算しない。

金属側の点電荷力は、\(\ell=1\) に正規化した別モジュール \texttt{ElectronMetal} である。\texttt{pullRate} は \(ke^2/\ell^3\) を単位に含んだ量であり、\texttt{pointForce} は \(k\) と \(e\) を別引数に取らない。
\begin{align}
\texttt{vecPhase}(w)&=\frac{1}{2\lVert w\rVert},\\
\texttt{pullRate}(x)&=\frac{8x^3}{\gamma_s^{\mathrm{eq}}(x)},\\
\texttt{pointForce}(Q,c,p)
&=-\bigl(Q\cdot\texttt{pullRate}(\texttt{vecPhase}(p-c))\bigr)\,(p-c).
\end{align}
\(r=\lVert p-c\rVert>0\) かつ \(\gamma_s^{\mathrm{eq}}>0\) かつ \(Q>0\) のとき、\(p\) に働くベクトルは \(c\) の向きを向く。この正規化での大きさは \(|Q|/( \gamma_s^{\mathrm{eq}} r^2)\) に一致する。確認は \(r=1/(2x)\) と \(8x^3\cdot r=4x^2\) と \(4x^2=1/r^2\) を使えばよい。\(r=0\) または \(\gamma_s^{\mathrm{eq}}=0\) の呼び出しは拒否する。\(\ell\neq 1\) の計算で \texttt{vecPhase} を流用しない。長さを戻すなら、位相を \(\ell/(2r)\) に取り、前置因子を単位に合わせて別関数にする。その関数は記録の \texttt{pointForce} と別名にする。

結晶の \texttt{siteForce} は \texttt{ElectronCrystal} にあり、有限集合 \(T\) 上のベクトル和
\[
\sum_{t\in T}(-q)\,\texttt{pointForce}(Q_t,\,p-\texttt{latticeVec}(b,h,t),\,x)
\]
である。\((-q)\) はベクトルへのスカラー倍であり、内積ではない。\(q\) は試験電荷、\(b\) は正規直交基底、\(h\) は格子間隔である。時間は進まない。無限格子和、遮蔽、凝集エネルギーは、この和の定義に入っていない。

\subsection{チャネルを一つの和にしない}

重力の \(N\) 体和と、電磁の \(N\) 体和は、別の関数として残す。呼び出し側が両方を欲しがるときは、二つのベクトルを受け取ったあとで足す。足し算そのものは第\ref{sec:contract}章の作業段に属する。Lean の \texttt{gravityLayerForce} は電磁項を含まず、\texttt{pointForce} の和は重力項を含まない。

\section{回帰検査}
\label{sec:checks}

この章は、\texttt{Regression.lean} が例として証明している数値範囲を、シミュレーターの単体テストへ写す。範囲そのものが新しい物理の主張ではない。テストは、主張の種類ごとにファイルを分ける。内在角のテストは \texttt{written\_at} だけを呼ぶ。層のテストは力の関数だけを呼ぶ。チャート比較のテストは、そのどちらも呼ばない。

\subsection{書かれた流れ}

次を機械的に検査する。

\begin{itemize}
\item \(\delta(t)=\delta_0+\nu t\) かつ \(\ddot\delta=0\)。
\item \(\sigma(t)\) の三次式と、\(\ddot\sigma=-2m\delta(t)\)。
\item \(\phi=(\sigma+\delta)/2\)、\(\theta=(\sigma-\delta)/2\)。
\item 数点の時刻で \(E(t)=E(0)=\frac12\nu\dot\sigma_0+\frac m2\delta_0^2\)。
\item \(p_\sigma(t)=\nu/2\) が一定。
\item サンプル \((m,\phi,\theta,\ddot\phi,\ddot\theta)=(1,1,0,-1,-1)\) は \texttt{PaperActualEL} を満たし、\(\ddot\delta=-m\delta\) を満たさない。
\item 同じサンプルの近傍 \((1,1,0,1,-1)\) は \texttt{PaperClaimedEL} を満たし、\texttt{PaperActualEL} を満たさない。
\item \(m\ge 0\) に対し \(\sqrt{2m}=m\) となるのは \(m=0,2\) だけである。この検査は振動子モジュールに置き、書かれた流れのテストから参照しない。
\end{itemize}

\subsection{層と力}

\begin{itemize}
\item \(\gamma_s^{\mathrm{eq}}(0)=1\)。
\item 導関数が \(-2\cosh x\sin x\) と、選んだ差分幅の中で一致する。
\item 見つけた第一正根が \(\pi/4\) と \(1\) のあいだにあり、その点で \(\gamma_s^{\mathrm{eq}}=0\)。
\item \(\kappa>0\)、\(r>0\)、\(\gamma<0\) で \texttt{gravityLayerTerm}\(>0\)。\(\gamma>0\) で項は負。
\item 二体の \texttt{gravityLayerPush} が反対称。
\item 斥力層で、外向き単位分離との内積が正。
\item \(N=3\) の和が、自分自身を項に含まない。
\item \(\gamma_G,\gamma_E,r\neq 0\) のとき、\texttt{pairChannelForce}\(=0\) と \(\kappa\gamma_E=\mathrm{coul}\,\gamma_G\) が同値。サンプルは半径を二種類に取り、零の有無が半径で変わらないことを見る。
\end{itemize}

\subsection{チャート比較}

\texttt{Regression} は次の区間を例として持つ。

\begin{itemize}
\item 水星の近日点余剰について \(42<\texttt{arcsecPerCentury}<44\)。
\item 太陽による光線の偏折について \(17/10<\texttt{deflectionArcsec}<9/5\)。
\end{itemize}

これらの数は、外部の軌道要素と SI スタンドインから出るチャート量である。内在角の閉形式の出力と比べない。テスト名に \texttt{chart} を入れる。

\subsection{核の長さが入る範囲}

散乱の符号が反転する読みでは、第一ノードの半径 \(r_1\) が
\[
\frac12<r_1<1
\]
（fm の入力として渡す）にあるとき、対応する長さ
\[
\texttt{lengthOfReversal}(r_1)=2\cdot\texttt{resonanceRoot1}\cdot r_1
\]
が
\[
\frac{14}{15}<\texttt{lengthOfReversal}(r_1)<\frac{19}{10}
\]
を満たす。これは、渡した \(r_1\) に対する長さの検査である。力を時間積分した出力ではない。\(\alpha_s\) をこの範囲から逆算しない。

強い力の結合定数を返す関数、核融合率を返す関数は、このリポジトリの記録に無い。テストもそれらを期待値にしない。

\section{作業段の仮定を残す}
\label{sec:contract}

地位は\dstopen{}である。Lean に積分器は無い。この章の一段は、実行してよい作業手順であり、論文の運動方程式ではない。仮定をログに残さない呼び出しは、ライブラリが拒否する。

\subsection{一段が読んでよいもの}

\begin{enumerate}
\item 位置はチャート \(E\) の \texttt{ChartBody.pos} だけである。
\item ペアの力は第\ref{sec:force}章の静的な和である。重力和と電磁和を両方使うなら、足し算をこの作業段で行い、その旨をログに書く。
\item 内在角は第\ref{sec:written}章の閉形式である。剛性 \(m\) は \texttt{WrittenSeed.stiffness} として外から渡す。層力の大きさから \(m\) を更新しない。
\item 位置を動かすときに加速度 \(\ddot x=F/\mu\) を使うなら、\(\mu\) は粒子ごとの慣性パラメータである。剛性 \(m\) でも符号なし質量 \(M\) でもない。
\end{enumerate}

\subsection{ログ}

各段の前に、次のレコードを残す。

\begin{verbatim}
struct StepContract:
    tag: "working-step"
    position_law: "frozen_accel_taylor"
    mu: map iota -> real
    stiffness_m: real
    killing_M: real
    used_gravity_sum: bool
    used_em_sum: bool
    note: "each mu_i, stiffness_m, and killing_M stay separate fields"
\end{verbatim}

\(\mu_i\)、\(m\)、\(M\) の数値が偶然一致しても、フィールドは分けたまま残す。一致を恒等式として簡約しない。

\subsection{一段の手順}

\begin{verbatim}
function working_step(bodies, seeds, mu, dt, inputs) -> reject or state:
    require step contract log
    for each body i:
        F_i = gravityLayerForce on the other bodies
        if em requested: F_i = F_i + em sum on the other bodies
        if any pair has r = 0 or gamma = 0: reject before writing
    for each body i:
        a = F_i / mu_i
        pos_i = pos_i + vel_i * dt + a * dt * dt / 2
        vel_i = vel_i + a * dt
    for each seed:
        sample written_at(seed, t + dt)
        build alpha, beta
        if not admissible: reject
    t = t + dt
\end{verbatim}

位置の差分は、段のあいだ加速度を一定とみなしたテイラー更新の例である。DST の定理ではない。別の段法を使うときは \texttt{position\_law} の文字列を変える。内在角の側は、段法を替えても閉形式のままにする。閉形式は錐の外でも式として定義されている。錐の検査は、シミュレーターが状態として保持するかどうかを決める門である。

\subsection{論文へ昇格させない}

この一段を、日本語論文の「書かれた作用」の節へ書き戻さない。Lean に \texttt{working\_step} という定義を足さない。多体系を閉じる式が足りない、という事実は、この仮定のログが引き受ける。

\section{数値法}
\label{sec:numerics}

\subsection{何を刻むか}

内在角は閉形式で評価する。ルンゲ--クッタで \(\phi,\theta\) を進め、閉形式と照合するテストは、実装の退行を見るためだけに置く。本番経路は閉形式である。

時刻方向に刻んでよいのは、チャート上の力和を位置へ渡す作業段だけである。段幅 \(\Delta t\) は入力である。\(\ell\) と \(m\) からクーラン数を導出しない。

\subsection{拒否}

次のいずれかが起きた段は、状態を更新せず拒否を返す。

\begin{itemize}
\item いずれかのペアで \(r=0\)。
\item いずれかのペアで、使用した \(\gamma_s^{\mathrm{eq}}\) または \(\gamma_G,\gamma_E\) が零（許容 \(\varepsilon_\gamma\) 以内）。
\item 更新後の \texttt{TorsionParams} が \texttt{IsAdmissibleContinuous} を満たさない。
\item 契約ログが無い。
\end{itemize}

拒否のあとで許容錐へ射影しない。最後に成功した状態を、呼び出し側が保持する。ライブラリは自動で段幅を半分にして再試行しない。再試行は呼び出し側のループが、新しい契約ログとともに行う。

\subsection{保存の検査}

検査してよい保存量は、一粒子・一軸のヤコビ積分 \(E\) と \(p_\sigma\) である。閉形式をサンプルした値が、初期データの定数から許容誤差を超えて離れたら、そのサンプルを失敗にする。

多体チャートの運動エネルギーと、層力の仕事の和が保存する、という主張は DST の定理としてログに書かない。作業段がエネルギー差分を監視するなら、監視値の名前を \texttt{chart\_kinetic\_plus\_work} とし、\(E\) と加算しない。

\subsection{特異点の近く}

\(\gamma_s^{\mathrm{eq}}\) が許容の外側で小さいとき、力は大きくなる。その大きさを頭打ちにする平滑化は、記録の式を別の式に取り替える操作である。平滑化を入れるなら \texttt{working} の別関数にし、\texttt{gravityLayerTerm} のテスト期待値を平滑化後の数にしない。

\section{離散双対は軌道積分ではない}
\label{sec:discrete}

地位は\dstrec{}である。ソースは \texttt{Algebra.Discrete} である。この章の手続きは時間ループに入らない。

\subsection{格子}

正整数 \(N\) に対し、ねじれの標本は \((\mathbb{Z}/N\mathbb{Z})^6\) である。
\begin{verbatim}
struct DiscreteTorsion:
    n: array[3] of mod N   # usual
    m: array[3] of mod N   # dual
\end{verbatim}
点数は \(N^6\) である。連続パラメータへの埋め込みは
\[
\alpha_a=\frac{2\pi}{N}(n_a\bmod N),\qquad
\beta_a=\frac{2\pi}{N}(m_a\bmod N)
\]
であり、識別子は \texttt{toTorsionParams} である。値は \([0,2\pi)\) に入る。許容かどうかは、埋め込んだあとの \texttt{IsAdmissibleContinuous} で判定する。格子点であること自体は許容を意味しない。

\subsection{トーラスは模型選択である}

循環指数は \(2\pi\) 周期である。双曲指数は \(2\pi\) 周期ではない。ブーストをこのトーラスへ載せることは、離散化の模型選択である。許容錐は、トーラスに載せたことから自動では従わない。列挙コマンドは、選択であることをヘルプ文字列に書く。

\subsection{コマンド}

\begin{verbatim}
command enumerate_admissible(N) -> list of DiscreteTorsion:
    for each of the N^6 points:
        p = to_torsion_params(point)
        if admissible(p): emit point
\end{verbatim}

出力は探索の入力である。各点の \(J\) と \(M\) を添えてよい。点を時刻順に並べ、速度として \texttt{working\_step} へ渡さない。\(N\) を段数と解釈しない。

埋め込みの成分は常に非負なので、離散点に限れば \texttt{Discrete.IsAdmissible} と、埋め込んだあとの \texttt{IsPrincipalBranch} は同値である。連続状態では同値ではない。列挙の検査は \texttt{IsAdmissible} を呼ぶ。

双対角だけを \((2\pi/N)\mathbb{Z}\) に制限し、通常ラピディティを実数のまま残す読みは、上の六次元格子とは別の模型である。同じコマンドにまとめない。

\section{モジュール境界と最小手順}
\label{sec:modules}

\subsection{分割}

推奨するソース分割は次のとおりである。矢印は、記録側が作業側を呼ばない向きだけを許す。

\begin{verbatim}
state
  torsion params, admissible test, J, M
written_flow
  written_at, paper energy, momenta
working_models
  runaway, oscillator, sourced jets, quartic
layer_force
  gamma_s_equal, gravityLayerTerm, gravityLayerForce
channels
  pairChannelForce, coulombForce, pointForce
step_contract
  working_step, contract log
checks
  one file per claim kind
discrete_enum
  enumerate_admissible
\end{verbatim}

\texttt{step\_contract} は \texttt{written\_flow} と \texttt{layer\_force} と、必要なとき \texttt{channels} を呼ぶ。\texttt{written\_flow} は \texttt{layer\_force} を呼ばない。\texttt{checks} は本番の段から呼ばれない。

\subsection{最小手順}

実装の順序は次に固定する。後の段に進む前に、前の段のテストが通る状態にする。

\begin{enumerate}
\item \texttt{state} と許容判定。ランダムな六つ組を許容と非許容に分け、境界 \(\alpha+\beta=\pi/2\) を一点ずつ確認する。
\item \texttt{written\_at}。第\ref{sec:checks}章の内在角テスト。
\item 二体の \texttt{gravityLayerTerm} と反対称な \texttt{gravityLayerPush}。
\item \(N\) 体の \texttt{gravityLayerForce}。自己項が無いこと。
\item 契約ログ付きの \texttt{working\_step} を、粒子二つ、一軸の種一つで一段だけ進める。拒否の四条件をそれぞれ一度は踏む。
\item チャート比較の区間テストを、内在角ソルバとリンクせずに置く。
\item 余力があれば \texttt{enumerate\_admissible} を小さい \(N\) で実行し、点数が \(N^6\) 以下であることと、許容点が錐の中にあることを見る。
\end{enumerate}

\subsection{名前}

公開関数の名前は、Lean の識別子と揃える。作業模型だけ接頭辞 \texttt{working\_} を付ける。チャート専用のテストには接頭辞 \texttt{chart\_} を付ける。同じ綴りで別の式を返す別名を作らない。

\section{埋めてはならない穴}
\label{sec:holes}

次の関数は、このライブラリの公開 API に存在しない。呼び出し側が研究のために実験するなら、別リポジトリの作業枝で行い、本書のモジュール名を使わない。

\begin{enumerate}
\item 多体の \(J\) を空間平均し、その平均をアインシュタイン方程式のソースとして返す関数。
\item 共有空間格子の上でマクスウェル方程式を時間発展させる関数。ファラデー係数を六つの生成子へ足す読みは、別の記録（\texttt{EMControl}）に属する。それを格子の PDE と同一視しない。
\item 層の半径から \(\alpha_s\) を返す関数。強い力の結合は導出されていない。
\item 粒子の生成、消滅、核融合、ベータ崩壊を一段の中で実行する関数。
\item チャート上の近日点余剰を、内在ローターの積分結果として返す関数。
\item \(G\)、\(\ell\)、\(m\)、\(a_0\)、微細構造定数を、許容錐または \(\gamma_s^{\mathrm{eq}}\) の零点から返す関数。
\item \(r\to 0\) の接触、合体、硬い反射。
\item 書かれた作用の \(E\) と、チャート上の力学的エネルギーを加算して「全エネルギー」と呼ぶ関数。
\item 量子重力稿のディラック方程式 \(\gamma^\mu\partial_\mu\Psi=-m\Psi\) を、多体状態の時間積分に使う関数。スピノルの記録は内在角の閉形式の代替ではない。
\end{enumerate}

穴が埋まったように見える数値は、第\ref{sec:contract}章のログが「作業」と書いた段の出力である。その出力を、Lean の定理番号の隣に置かない。

\section{付録}
\label{sec:appendix}

\subsection{式と Lean 名}

\begin{description}
\item[六ラピディティ] \texttt{TorsionParams}（\texttt{Algebra.Operations}）
\item[許容錐] \texttt{IsAdmissibleContinuous}（\texttt{Algebra.Admissible}）
\item[\(J\)、\(J_{\mathrm{norm}}\)、\(M\)] \texttt{J}, \texttt{JNormalized}, \texttt{mass}（\texttt{Algebra.Invariant}）
\item[書かれた密度] \texttt{paperLagrangian}（\texttt{DualRotorDynamics}）
\item[書かれた EL] \texttt{PaperActualEL}（\texttt{DualRotorDynamics}）
\item[\(p_\sigma\)、\(p_\delta\)、\(E\)] \texttt{paperMomentumSigma}, \texttt{paperMomentumDelta}, \texttt{paperEnergy}（\texttt{DualRotorDynamics}）
\item[\(\delta(t)\)、\(\sigma(t)\)] \texttt{writtenDelta}, \texttt{writtenSigma}（\texttt{DualRotorFlow}）
\item[暴走、振動子] \texttt{runawayDelta}, \texttt{harmonicMismatch}（\texttt{DualRotorFlow}）
\item[駆動 \(u,v\)] \texttt{PaperSourcedEL}, \texttt{PaperBothSourcedEL}（\texttt{DualRotorDynamics}）
\item[任意傾き] \texttt{OppositePotentialEL}, \texttt{SameSignPotentialEL}（\texttt{DualRotorRigidity}）
\item[四次井戸] \texttt{quarticPotential}, \texttt{quarticSlope}（\texttt{DualRotorRigidity}）
\item[等スケール因子] \texttt{gammaSEqual}（\texttt{TorsionalLayer}）
\item[第一正根] \texttt{resonanceRoot1}（\texttt{ElectronShell}）
\item[層半径] \texttt{layerRadius}（\texttt{TorsionalLayer}）
\item[重力項、和] \texttt{gravityLayerTerm}, \texttt{gravityLayerForce}（\texttt{LayerForce}）
\item[二チャネル] \texttt{pairChannelForce}（\texttt{BothChannels}）
\item[クーロンスカラー] \texttt{coulombForce}（\texttt{ElectronOrbit}）
\item[点電荷] \texttt{pointForce}（\texttt{ElectronMetal}）
\item[関係位置] \texttt{S3}, \texttt{relative}（\texttt{RelationalPosition}）
\item[離散標本] \texttt{DiscreteTorsion}（\texttt{Algebra.Discrete}）
\item[剛性ポテンシャル] \texttt{stiffnessPotential}（\texttt{RestMass}）
\item[外部有理数] \texttt{SI}（\texttt{Gravity.SI}）
\end{description}

\subsection{章と論文}

\begin{description}
\item[状態、許容錐、\(J\) と \(M\)] 主論文の数学的基礎と運動学。
\item[書かれた作用、近傍模型、駆動 \(u\)、任意の \(V\)] 動態方程式の論文。
\item[層、力、核の長さの範囲、電子の点電荷] 主論文の強場域、スケール不変層、クーロン層。
\item[チャートの近日点と偏折] 主論文の弱場・ニュートン極限と、Lean の \texttt{Mercury}, \texttt{SolarDeflection}。
\item[離散列挙] 離散双対時空の草稿と \texttt{Algebra.Discrete}。
\item[ディラック方程式を時間積分に使わない] 量子重力稿のクリフォード平方の節。
\end{description}

量子重力稿の場の方程式は、双対スピノルのチャート上の読みである。多体シミュレーターの状態ベクトルへ \(\Psi\) を追加しない。

\subsection{照合}

本書に書いた閉形式、ヤコビ積分、層力、二チャネルの零条件、許容錐、離散埋め込みは、上表の Lean 定義および動態方程式の日本語稿と一致している。一致しなかった式は採用していない。作業段の差分法は Lean に無く、仮定のログの外へ出さない。


\end{document}
